\documentclass[a4paper,12pt]{article}
\usepackage[utf8]{inputenc}
\usepackage[russian]{babel}
\usepackage{amsmath, amssymb}

\title{Лекция 5 \\ Решетчатые диаграммы линейных кодов}
\author{}
\date{}

\begin{document}

\maketitle

\section{Признак минимальности решетки}

Решетчатая диаграмма (треллис) линейного кода представляет собой ориентированный граф,
развернутый во времени, который отражает все кодовые слова. Состояния на каждом ярусе
(временном слое) соответствуют некоторым подпространствам, а ребра --- переходам,
порождающим символы кодового слова.

Минимальная решетка --- это решетка с наименьшим возможным числом состояний на каждом ярусе.

\textbf{Основной признак минимальности:} на каждом ярусе все пути, имеющие общее прошлое
или общее будущее, проходят через один узел.

Иначе говоря, два различных состояния на одном ярусе не могут быть объединены, поскольку
множества путей, продолжающихся из них в будущее (или приходящих из прошлого), различны.

Для любого линейного кода существует единственная (с точностью до изоморфизма)
минимальная решетка.

\section{Построение решетки по порождающей матрице}

Пусть задан линейный $(n,k)$-код $C$ над полем $\mathbb{F}_2$ с порождающей матрицей
$G$ размера $k \times n$.

Для построения решетки необходимо привести строки матрицы к минимальной спэновой форме (МСФ).

Для каждой $i$-й строки определяются:
\begin{itemize}
    \item $b_i$ --- индекс первого ненулевого элемента (начало спэна),
    \item $e_i$ --- индекс последнего ненулевого элемента (конец спэна).
\end{itemize}

Спэн строки --- это интервал $[b_i, e_i]$.

После упорядочивания строк по возрастанию $e_i$ можно построить решетку с ярусами
$j = 0,1,\dots,n$.

Состояние на ярусе $j$:
\[
S_j = \langle g_i \mid b_i \le j < e_i \rangle
\]

Начальное и конечное состояния:
\[
S_0 = S_n = \{0\}
\]

Переход между состояниями определяется выбором информационного символа.

\section{Пример}

Рассмотрим матрицу:
\[
G =
\begin{pmatrix}
1 & 0 & 1 & 1 & 0 & 1 \\
1 & 0 & 1 & 0 & 1 & 0 \\
1 & 1 & 0 & 1 & 0 & 0
\end{pmatrix}
\]

Здесь $k = 3$, $n = 6$, скорость кода:
\[
R = \frac{k}{n} = \frac{1}{2}
\]

Все кодовые слова $c = uG$, где $u \in \mathbb{F}_2^3$:

\[
\begin{array}{ccc|c}
m_1 & m_2 & m_3 & \text{код} \\
\hline
0 & 0 & 0 & 000000 \\
0 & 0 & 1 & 000111 \\
0 & 1 & 0 & 011110 \\
0 & 1 & 1 & 011001 \\
1 & 0 & 0 & 110100 \\
1 & 0 & 1 & 110011 \\
1 & 1 & 0 & 101010 \\
1 & 1 & 1 & 101101 \\
\end{array}
\]

Минимальное расстояние:
\[
d_{\min} = 3
\]

\subsection{Спэны строк}

\begin{itemize}
    \item Строка 1: позиции $1,3,4,6$ $\Rightarrow b_1 = 1,\ e_1 = 6$
    \item Строка 2: позиции $1,3,5$ $\Rightarrow b_2 = 1,\ e_2 = 5$
    \item Строка 3: позиции $1,2,4$ $\Rightarrow b_3 = 1,\ e_3 = 4$
\end{itemize}

После сортировки: $g_3, g_2, g_1$.

\subsection{Состояния по ярусам}

\begin{itemize}
    \item $j=0$: $S_0 = \{0\}$
    \item $j=1,2,3$: активны все строки $\Rightarrow S_j = \mathbb{F}_2^3$
    \item $j=4$: $S_4 = \langle g_2, g_1 \rangle$
    \item $j=5$: $S_5 = \langle g_1 \rangle$
    \item $j=6$: $S_6 = \{0\}$
\end{itemize}

Профиль решетки:
\[
1, 8, 8, 8, 4, 2, 1
\]

\section{Теорема о минимальности}

\textbf{Теорема.} Если порождающая матрица приведена к МСФ, то построенная решетка является минимальной.

\textbf{Доказательство.}

Состояние на ярусе $j$ определяется линейной комбинацией активных строк.
Если два состояния различны, то существует строка, коэффициенты при которой различаются.

Поскольку эта строка активна до $e_i$, дальнейшее поведение различается,
следовательно множества путей различны.

Значит состояния нельзя объединить, и решетка минимальна.

\end{document}